\section{Introduction}\label{sec:intro}

Imagine a scientist who regresses a response on $p$ covariates observed over $n$ units, with $p$ comparable to, or larger than, $n$: several thousand gene expressions, hundreds of firm or asset characteristics, brain-imaging summaries, text embeddings. The covariates were not randomized, so it is plausible that a handful of unmeasured common causes move many covariates and the response simultaneously, a situation we call dense latent confounding. Three questions follow. First, is the confounding biasing my estimates, and by how much? Second, can I tell from the observed data alone that it is there? Third, if I adjust spectrally, for instance by trimming leading principal components, does the adjustment help me or hurt me? Common practice answers the second question by inspecting the scree plot of the sample eigenvalues. A recurring message of this paper is that this practice fails quantitatively in both directions: scree-type evidence fires on ordinary design structure whether or not confounding is present, and it stays silent exactly where silence is dangerous.

We study the standard latent-factor model for one observational dataset,
\[
X = \Lambda f + u, \qquad Y = \beta' X + \gamma' f + \varepsilon ,
\]
where $X$ collects the $p$ covariates of a generic unit, $Y$ is the scalar response, $f$ stacks $r$ hidden factors with $r$ fixed, $u$ is covariate noise with variance $\sigma_u^2$, and $\varepsilon$ is response noise; only $(Y,X)$ is observed. Here $\Lambda$ is a dense $p \times r$ loading matrix, $\beta \in \mathrm{R}^p$ is the structural coefficient vector the scientist wants, and $\gamma \in \mathrm{R}^r$ is the confounder-response link. The aspect ratio is $c = p/n$, and $\Sigma_X = \sigma_u^2 I + \Lambda\Lambda'$ is the population covariance of $X$. The spike strengths $l_1 \ge \dots \ge l_r \ge 0$ are the eigenvalues of $\Lambda\Lambda'/\sigma_u^2$, so the population eigenvalues of $\Sigma_X/\sigma_u^2$ exceed one exactly by the $l_j$. Informally, $l_j$ measures how far the $j$-th hidden factor stands above the noise floor, and $\sqrt{c}$ is the visibility threshold inherited from random matrix theory: a factor with $l_j < \sqrt{c}$ leaves no outlier eigenvalue, while one with $l_j > \sqrt{c}$ produces a separated spike. Two objects summarize the confounding itself: the bias direction $b = \Lambda\gamma$, the $p$-vector through which hidden factors enter the conditional mean of $Y$, and $g = \|\gamma\|_2$, the overall confounding strength; the alignment angle $\theta$ between $\gamma$ and the dominant factor direction determines where the confounding mass rides.

The geometry behind our results is a decoupling. Visibility depends only on the spectral profile $(l_j, c)$, because eigenvalue outliers do not see $\gamma$. Harmfulness depends on $b = \Lambda\gamma$. Since these are different functions of the parameters, there are regions of $(l, c, g, \theta)$ where the confounding is invisible to any eigenvalue inspection yet biases ordinary least squares by a large constant, and regions where a spectacular spike is visible yet essentially harmless. A practitioner equipped only with a scree plot cannot tell these worlds apart.

Stated in one sentence, this paper contributes exact spectral theory for hidden dense confounding: a validated capture law for the bias confounding induces, a calibrated detection test whose power frontier is predictable to within a factor of about one, a proof-backed map of what cannot be detected and by which classes of statistics, and a trim-then-regress recipe for treatment effects, with the failure of tuned soft-trim estimators characterized rather than hidden. In detail:

\begin{enumerate}
\item \emph{An exact capture law for confounding bias.} For ordinary least squares at aspect ratios $c > 1$, the confounding-induced bias along the $j$-th spike direction is capped exactly by $\mathrm{cap}_j = (1+l_j)/(c+l_j)$, and at $c \le 1$ an exact identity pins the bias down completely. The law is proved elementarily at $r=1$ (a corrected channel assembly plus a Haar-probe symmetry argument, with no local-law input) and reduced at fixed $r$ through Woodbury identities and vanishing lemmas. A companion ridge formula interpolates the cap continuously in the regularization parameter and collapses, as the regularization vanishes, to the min-norm OLS law; a trimmed-$\tau$ law extends the picture to scalar treatment effects, predicting the OLS attenuation denominator excess $\Lambda_D = 1.558$ against a measured $1.5548 \pm 0.072$. Simulation overlays confirm the theory: 91.7\% of the main-grid cells fall within 10\% of the predicted bias at $n=2000$ and every cell does so at the neighboring tiers, with ridge deviations at median 0.0\%.
\item \emph{A calibrated detection alarm with a predictable power frontier.} Our statistic S2, a calibrated maximum of standardized spike coordinates, ships with a finite-sample null law: its Monte Carlo 95th percentile exceeds the naive Gaussian-maximum scale by a factor of about 1.7, which a noise-aware recalibration captures (pooled chi-square $p = 0.41$; 96.8\% of null cells inside the $|z| \le 2$ band). Its predicted power frontier matches simulation: the median ratio of empirical to predicted 80\%-power strengths is 1.02, with stratum ratios 0.93, 1.25, and 1.02, all within the predeclared factor 1.5.
\item \emph{A map of what cannot be detected.} We prove a population detachment boundary for augmented-spectrum alarms (Proposition~\ref{prop:detach}): for subcritical dense confounding under response standardization no eigenvalue detaches from the bulk at any strength, the low-edge wake criterion fails at every studied aspect ratio, and the limiting bulk spectra under null and alternative coincide. A stronger universal impossibility in the style of Onatski--Moreira--Hallin (OMH) \citep{onatski2013asymptotic} was planned, pre-registered, and destroyed by our own falsifiers: the true augmented top eigenvalue separates consistently below the boundary through finite-sample fluctuations that leave the bulk laws untouched, with geometry-dependent direction, and a pipeline erratum behind the frozen S1 statistic explains part of how the wrong picture survived. What ships is layered: blindness is trivial for purely covariate-based alarms, proved for the spike-coordinate calibrated alarm we ship (power ceiling confirmed in nine of nine blind strata), and false for the augmented top eigenvalue. Complementing this map, an exact visibility law for quadratic probes exhibits satiation and places the born-blind crossing of our scaling calibration exactly at $c=1$; learning-based probes confirm operational invisibility at $c=0.8$ through $g\approx0.8$ while detecting readily at $c=0.2$ (AUC up to 1.0), so probe-level invisibility is claimed only in intermediate-to-high-$c$ regions.
\item \emph{A characterized negative result for soft tuning, and a recipe that survives.} Empirical-Bayes-tuned soft spectral trimming, the natural adaptive member of the deconfounding family, collapses to OLS or loses to hard trimming in every harmful cell of our grids: its no-regret share (the share of cells with mean-bias norm within $1.05\times$ the best baseline's) is 4.1\% against a predeclared requirement of 95\%. A collapse lemma explains part of the failure mechanically (the empirical-Bayes ratio vanishes under our primary ledger, driving the soft weights to one, so the tuned estimator reproduces OLS identically in 31 of the 97 harmful cells), and a dominance theorem explains why Onatski hard trimming wins outright (best baseline in 89 of 97 harmful cells). What survives is trim-then-regress: treatment error stays flat in the aspect ratio under hard trimming while OLS error inflates about six-fold, and on real geometry the trimmed estimate is 0.173 against OLS 0.421 at a true treatment effect of 1 (ratio 2.43, correct sign in 100\% of replicates).
\item \emph{Semi-synthetic benchmarks and a released package.} We inject known dense confounding into three real covariate geometries (AddNeuroMed gene expression, a featurized IHDP covariate set, and the k401ksubs wealth panel) across six configurations. Before injection, every configuration is massively outlier-positive and the scree alarm rejects always (Tracy-Widom statistics up to 48,431), including under matched nulls and permutation, confirming that spectrum inspection carries no information about confounding. After injection, alarm power straddles each family's own predicted frontier (power between 0.14 and 0.21 at half the predicted strength, and between 0.975 and 1.00 at twice it) while size stays between 0.000 and 0.003 under every negative control across roughly 3,000 control evaluations; incumbent confounding-strength diagnostics track injected strength monotonically but turn anti-conservative out of sample (rejection rates 0.07--0.92 and 0.61--1.00). The pipeline ships as the \texttt{confounderalarm} package (v0.1.0), which returns a verdict, a p-value, estimates of $(r, l, c)$, frontier placement, a blind-region certificate, and a recommended Onatski-trim adjustment; every benchmark cell replicates to machine precision (largest statistic shift $7.7 \times 10^{-13}$) on independent hardware.
\end{enumerate}

The remainder of the paper is organized as follows. Section 2 fixes the models, normalizations, estimands, and the assumption ledger, including the identification subtlety that only a projection of $\beta$ is determined by second moments. Section 3 states and proves the capture law, its ridge interpolation, and the trimmed-$\tau$ law for treatment effects. Section 4 constructs the calibrated alarm and compares measured power against the predicted frontier. Section 5 maps what cannot be detected: a proved detachment boundary, a falsified universal impossibility reported as an instructive negative result together with its pipeline erratum, the class-scoped blindness of the alarm we ship, and the quadratic-probe visibility boundary. Section 6 reports the estimation negative-result analysis: why oracle-tuned soft trimming fails, why hard trimming dominates, and what survives. Section 7 validates the toolkit on the semi-synthetic benchmarks and describes the package. Section 8 collects limitations, and Section 9 discusses how our tests complement, rather than replace, sensitivity-analysis practice.

\subsection{Related work}\label{sec:related}

Closest to us is the spectral deconfounding family. \citet{cevid2020spectral} showed that dense hidden confounding destroys consistency of sparse high-dimensional regression and that suitable spectral transformations of the design restore it. The idea has been extended to sparse additive models by \citet{scheidegger2024spectral}, carried into random forests through a spectral loss by \citet{ulmer2025sdforest}, and combined with boosting and an empirical-Bayes tuning of the spectral loss by \citet{nava2026sdboost}. These are estimation-method papers: they deliver rates and consistency guarantees, often with arbitrary or heuristic default tuning, and none of them contains a bias phase diagram, a detectability or removability frontier, a calibrated test with controlled size, or an impossibility statement. Our optimality verdict is positioned squarely against this family, in particular against the empirical-Bayes tuning of \citet{nava2026sdboost}, and Section 6 shows that tuned soft weights lose comprehensively to hard trimming on harmful cells.

A second line quantifies confounding strength. \citet{janzing2018detecting} estimate the degree of confounding from an independence-of-mechanisms principle via spectral concentration arguments, and \citet{rendsburg2024consistent} establish consistency of a random-matrix-corrected version of that estimator. This line stops at description: it ships no exact null distribution, no law for the induced bias, no statement about when confounding matters for a specific causal estimand, and no guidance about removal.

A third neighborhood is high-dimensional endogeneity and latent-factor inference. \citet{fan2014endogeneity} construct factor-adjusted projected statistics for testing under strong factors, \citet{guo2022ddll} provide debiased inference for sparse coefficients under hidden confounding, and \citet{wang2019blessings} develop the multiple-causes identification argument. Their targets are significance, interval estimation, and identification; none of them analyzes the spectral phase behavior of bias, detectability, or removability, which is the object here.

Methodologically we compose the random matrix toolbox into causal-estimand functionals: the Marchenko-Pastur law \citep{marchenko1967distribution}, the Baik-Ben Arous-Pech\'e transition \citep{bbp2005phase}, Tracy-Widom flank distributions \citep{johnstone2001distribution}, outlier locations and eigenvector overlaps \citep{benaychgeorges2011eigenvalues,benaychgeorges2012singular}, sharp power analysis of spectral tests near the transition \citep{onatski2013asymptotic}, consistent estimation of the number of spikes \citep{onatski2010determining}, exact ridge risk \citep{dobriban2018high,hastie2022surprises}, and anisotropic local laws \citep{knowles2017anisotropic}. Our contribution is not any single one of these laws but their assembly into bias caps for OLS, ridge, and spectral trims, and into power frontiers for confounding alarms.

To our knowledge, no prior work provides the combination delivered here: a bias phase diagram, a detectability frontier with a structural removability verdict, a calibrated alarm, and an impossibility map stating which classes of statistics are blind where. Targeted vocabulary searches support the gap claim: arXiv API queries executed on 2026-08-23 returned zero hits for ``deconfounding'' AND ``phase transition'', zero for ``hidden confounding'' AND ``random matrix'', zero each for ``confounding'' paired with ``Benaych-Georges'' or ``Tracy-Widom'', and zero for ``confounder detection'' AND ``spike'', while ``spectral deconfounding'' returns exactly the four-paper family above; a forward-citation screen of the family's records (65 unique works merged across two APIs) surfaced no phase-transition or frontier analysis either.
